\section{机器人基础模型}

\subsection{从专用策略到基础模型}

\begin{figure}[H]
\centering
\includegraphics[width=0.95\textwidth]{slides-images/slide-059.jpg}
\caption{机器人基础模型的发展时间线：RT-1 $\rightarrow$ RT-2 $\rightarrow$ RT-X $\rightarrow$ OpenVLA $\rightarrow$ $\pi_0$ 等\protect\footnotemark}
\end{figure}
\footnotetext{Slides 第60页。}

\begin{importantbox}{机器人基础模型的定义}
机器人基础模型是一种广泛泛化的策略，其输出的动作可能不是最优的，但在给定任何合理的观测和语言指令条件下，\textbf{总是能生成平滑、合理、可在真实世界执行的动作轨迹}。这类似于大语言模型——回答可能不完美，但总是``有道理''的。
\end{importantbox}

这些模型有多种名称（Vision-Language-Action Models / VLAs、大行为模型 / Large Behavior Models），但本质相同：从观测和任务指令映射到可广泛泛化的动作输出。

\subsection{$\pi_0$：一个具体案例}

\begin{figure}[H]
\centering
\includegraphics[width=0.95\textwidth]{slides-images/slide-062.jpg}
\caption{$\pi_0$ 框架：大规模多源数据预训练 $+$ 视觉语言模型初始化 $+$ 任务特定后训练\protect\footnotemark}
\end{figure}
\footnotetext{Slides 第63页。}

$\pi_0$（Physical Intelligence 于 2024 年 10 月发布）的设计包含三个关键阶段：

\begin{enumerate}
    \item \textbf{数据聚合}：汇集学术界和自采集的大量多机器人、多任务数据
    \item \textbf{预训练}：从预训练的视觉语言模型出发，用动作预测和 VQA 双重目标共同微调——既保留语义知识，又学会预测动作
    \item \textbf{后训练（Post-training）}：收集特定任务数据，对基座模型进行微调以达到高性能
\end{enumerate}

\begin{figure}[H]
\centering
\includegraphics[width=0.95\textwidth]{slides-images/slide-063.jpg}
\caption{$\pi_0$ 的三级评估：零样本基座模型、后训练的分布内任务、后训练的新任务\protect\footnotemark}
\end{figure}
\footnotetext{Slides 第64页。}

\begin{knowledgebox}{为什么基础模型路线可行？}
机器人基础模型的成功来自三个因素：
\begin{enumerate}
    \item \textbf{视觉语言模型的语义知识迁移}：预训练 VLM 已理解``抽屉''、``把手''等概念，这些知识可直接用于机器人任务
    \item \textbf{大规模多源数据}：跨机器人、跨任务的数据带来泛化能力
    \item \textbf{扩散模型等表达力强的策略类}：能处理多模态动作分布
\end{enumerate}
\end{knowledgebox}

\subsection{当前挑战}

\begin{figure}[H]
\centering
\includegraphics[width=0.95\textwidth]{slides-images/slide-068.jpg}
\caption{机器人学习面临的核心挑战：数据规模不足、泛化鸿沟、评估标准化缺失\protect\footnotemark}
\end{figure}
\footnotetext{Slides 第69页。}

尽管进展迅速，机器人基础模型仍面临重大挑战：

\begin{itemize}
    \item \textbf{数据不足}：机器人数据的规模远不及语言/视觉数据。数据收集依赖人类遥操作，速度比人类直接操作更慢
    \item \textbf{泛化不够}：在``野外''环境（非实验室设置）中的泛化能力仍不足
    \item \textbf{评估困难}：缺乏标准化的大规模评估基准，难以客观比较不同方法的进展
    \item \textbf{执行效率}：当前机器人执行任务远慢于人类
\end{itemize}

\subsection{本章小结}

机器人基础模型借鉴大语言模型的成功经验，通过大规模多源数据预训练和任务特定后训练，追求广泛泛化的机器人策略。$\pi_0$ 等工作展示了令人鼓舞的进展，但数据规模、泛化能力和评估标准仍是关键瓶颈。

%% ========== 总结与延伸 ==========

